\documentclass[11pt,a4paper]{article}
\usepackage[T1]{fontenc}
\usepackage[utf8]{inputenc}
\usepackage[brazil]{babel}
\usepackage{geometry}
\usepackage{xcolor}
\usepackage{enumitem}
\usepackage{titlesec}
\usepackage{needspace}
\usepackage{hyperref}

% ------------------------------------------------------------------
% Geometria da pagina (layout confortavel de duas paginas)
% ------------------------------------------------------------------
\geometry{top=0.55in, bottom=0.5in, left=0.65in, right=0.65in}
\pagestyle{empty}
\setlength{\parindent}{0pt}
\setlength{\parskip}{0pt}

% ------------------------------------------------------------------
% Cores, links e estilo das secoes
% ------------------------------------------------------------------
\definecolor{primary}{RGB}{31, 73, 125}
\definecolor{muted}{RGB}{90, 90, 90}
\hypersetup{colorlinks=true, urlcolor=primary, linkcolor=primary}

\titleformat{\section}{\large\bfseries\color{primary}\uppercase}{}{0pt}{}[\titlerule]
\titlespacing{\section}{0pt}{9pt}{5pt}

% Cabecalho de empresa (nome + local) -- mantido junto com o cargo/bullets seguintes
\newcommand{\company}[2]{%
  \needspace{5\baselineskip}%
  \textbf{#1} \hfill \textcolor{muted}{#2}\\[3pt]
}
% Cargo dentro de uma empresa (titulo + datas)
\newcommand{\subrole}[2]{%
  \needspace{4\baselineskip}%
  \textit{#1} \hfill \textcolor{muted}{#2}\\[3pt]
}
% Cargo em linha unica (titulo + datas + local) -- mantido junto com seus bullets
\newcommand{\role}[4]{%
  \needspace{5\baselineskip}%
  \textbf{#1} \hfill \textbf{#3} \\[2pt]
  \textit{#2} \hfill \textcolor{muted}{#4}\\[4pt]
}
% Entrada de projeto
\newcommand{\project}[3]{%
  \textbf{#1} \,\textcolor{muted}{\small[#2]} \, --- \, #3 \\[3pt]
}

\begin{document}

% ------------------------------------------------------------------
% Cabecalho
% ------------------------------------------------------------------
\begin{center}
    {\Huge\bfseries Za\'u J\'ulio} \\[5pt]
    \small Tech Lead \,|\, Engenheiro de Software Full Stack \,|\, Microsservi\c{c}os \& Design de Sistemas \,|\, IA / ML \\[4pt]
    \footnotesize
    \href{mailto:zaujulio.dev@gmail.com}{zaujulio.dev@gmail.com} \,|\, +55 84 99865-1868 \,|\, Jo\~ao Pessoa, PB, Brasil \\[2pt]
    \href{https://linkedin.com/in/zaujulio}{linkedin.com/in/zaujulio} \,|\,
    \href{https://github.com/ZauJulio}{github.com/ZauJulio} \,|\,
    \href{https://profile.codersrank.io/user/zaujulio}{codersrank.io/@zaujulio} \,|\,
    \href{https://zaujulio.com}{Portf\'olio: zaujulio.com}
\end{center}

% ------------------------------------------------------------------
% Resumo
% ------------------------------------------------------------------
\section{Perfil Profissional}
Tech Lead e Engenheiro de Software Full Stack m\~ao na massa, com mais de 5 anos construindo sistemas escal\'aveis e de alta concorr\^encia ao longo de todo o ciclo de vida do software. Profundo dom\'inio do ecossistema \textbf{TypeScript / Node.js} (\textbf{GraphQL}, Express, NestJS) e de \textbf{Python} para \textbf{IA, Machine Learning, Deep Learning, NLP e s\'eries temporais}. Projeta \textbf{microsservi\c{c}os} orientados a eventos com filas de processamento (BullMQ, RabbitMQ), modela dados relacionais e NoSQL (PostgreSQL, MySQL/MariaDB, MongoDB) e \'e respons\'avel pela entrega de ponta a ponta---da modelagem Entidade-Relacionamento e servi\c{c}os \textbf{C\# / .NET} a CI/CD, infraestrutura Linux e otimiza\c{c}\~ao de custos em nuvem. Entre os \textbf{melhores do mundo em GraphQL no CodersRank}. Come\c{c}ou a programar antes da onda de IA e hoje \'e fluente em desenvolvimento assistido por IA, liderando times sem deixar de entregar c\'odigo em produ\c{c}\~ao diariamente.

% ------------------------------------------------------------------
% Competencias Tecnicas (no topo para leitura rapida de recrutador / ATS)
% ------------------------------------------------------------------
\section{Compet\^encias T\'ecnicas}
\begin{itemize}[leftmargin=14pt, itemsep=2pt, topsep=0pt]
    \item \textbf{Linguagens de Programa\c{c}\~ao:} TypeScript / JavaScript, Python, C\#, Java, C, Bash, SQL.
    \item \textbf{Backend \& Microsservi\c{c}os:} Node.js, Express, NestJS, GraphQL, .NET / .NET Core, FastAPI, Django, REST, WebSockets / Socket.io, BullMQ, RabbitMQ.
    \item \textbf{Frontend:} React, Vite, Next.js, Redux, Tailwind CSS, Styled Components, HTML / CSS, SSR, Core Web Vitals.
    \item \textbf{Bancos de Dados:} PostgreSQL, MySQL / MariaDB, MongoDB, Redis, Data Lakes, modelagem ER.
    \item \textbf{IA / Dados:} Machine Learning, Deep Learning, NLP, S\'eries Temporais, Pandas, NumPy, scikit-learn, PyTorch, Matplotlib, Jupyter, R.
    \item \textbf{Infra / DevOps:} Docker, Linux, Shell script, CI/CD, AWS, GCP, VPC, CDN, DevSecOps, Monitoramento / Observabilidade.
    \item \textbf{Pr\'aticas:} Design de sistemas, arquitetura de microsservi\c{c}os, SDLC, code review, lideran\c{c}a t\'ecnica, desenvolvimento assistido por IA (Claude, OpenCode).
\end{itemize}

% ------------------------------------------------------------------
% Experiencia
% ------------------------------------------------------------------
\section{Experi\^encia}

\company{Bitwise Technology}{Sorocaba, SP (Remoto)}
\subrole{Tech Lead}{Ago 2025 -- Atual}
\begin{itemize}[leftmargin=14pt, itemsep=2pt, topsep=0pt]
    \item Defino o roadmap de engenharia e os padr\~oes t\'ecnicos de microsservi\c{c}os TypeScript e schemas GraphQL, sendo respons\'avel pelo design de sistemas de ponta a ponta (da modelagem ER ao deploy) sem deixar de ser m\~ao na massa, entregando c\'odigo em produ\c{c}\~ao diariamente.
    \item Conduzi otimiza\c{c}\~oes de infraestrutura direcionadas que reduziram os tempos de build de CI/CD em 75\% (de 20 min para menos de 5) e o custo mensal na AWS em \textasciitilde80\% (de US\$120 para US\$25), al\'em de redimensionar servi\c{c}os Node.js em produ\c{c}\~ao para metade do consumo de CPU e mem\'oria sem impacto na vaz\~ao.
    \item Atuo na mentoria de times multidisciplinares e conduzo code reviews, promovendo cultura de qualidade, escalabilidade e observabilidade.
\end{itemize}
\vspace{3pt}

\subrole{Engenheiro de Software Full Stack}{Mar 2023 -- Ago 2025}
\begin{itemize}[leftmargin=14pt, itemsep=2pt, topsep=0pt]
    \item Entreguei backends de alta disponibilidade em GraphQL, Node.js, C\#/.NET e MongoDB, modelando regras de neg\'ocio complexas e delegando processamento pesado a filas BullMQ e RabbitMQ.
    \item Constru\'i frontends React/Vite/Next.js r\'apidos e otimizados para SSR e Core Web Vitals, al\'em dos pipelines de CI/CD e do provisionamento em nuvem que encurtaram os ciclos de release.
\end{itemize}
\vspace{3pt}

\role{Commcepta}{Desenvolvedor Full Stack}{Dez 2021 -- Dez 2022}{Curitiba, PR (Remoto)}
\begin{itemize}[leftmargin=14pt, itemsep=2pt, topsep=0pt]
    \item Otimizei consultas GraphQL sobre um Data Lake de larga escala, reduzindo a lat\^encia de buscas de \textasciitilde30s para menos de 6s (at\'e 50\% mais r\'apido) e viabilizando an\'alises quase em tempo real sobre Big Data.
    \item Desenvolvi pipelines em Python que unificaram m\'ultiplas fontes de dados para a tomada de decis\~ao corporativa, e refor\c{c}ei o portal de uma multinacional em acessibilidade (a11y) e internacionaliza\c{c}\~ao (i18n).
\end{itemize}
\vspace{3pt}

\role{Byte Serid\'o Jr -- Empresa J\'unior (UFRN)}{Diretor de Projetos / Desenvolvedor Front-end}{Jun 2021 -- Nov 2022}{Caic\'o, RN}
\begin{itemize}[leftmargin=14pt, itemsep=2pt, topsep=0pt]
    \item Liderei o planejamento e a entrega de m\'ultiplos projetos de software na empresa j\'unior da universidade, evoluindo de desenvolvedor front-end a Diretor de Projetos.
    \item Mentorei e tutorei novos desenvolvedores, integrando estudantes em fundamentos de programa\c{c}\~ao, boas pr\'aticas de engenharia e fluxos de projeto.
\end{itemize}

% ------------------------------------------------------------------
% Pesquisa & Publicacao
% ------------------------------------------------------------------
\section{Pesquisa \& Publica\c{c}\~ao}

\role{LabEPI / UFRN}{Pesquisador de IA \& Machine Learning (Inicia\c{c}\~ao Cient\'ifica)}{Ago 2019 -- Set 2020}{Caic\'o, RN}
\begin{itemize}[leftmargin=14pt, itemsep=2pt, topsep=0pt]
    \item Projetei um classificador h\'ibrido combinando redes neurais \textbf{Self-Organizing Maps (SOM)} com Regress\~ao Polinomial para detec\c{c}\~ao n\~ao supervisionada de anomalias em dados de energia minuto a minuto, superando a acur\'acia de um baseline K-Means e reduzindo falsos positivos.
    \item Otimizei a gera\c{c}\~ao e reutiliza\c{c}\~ao de modelos de IA, reduzindo o consumo de mem\'oria de 2 GB para \textasciitilde200 MB e o tempo de gera\c{c}\~ao das redes SOM de 10--15 minutos para menos de 2 minutos.
    \item Autor de artigo revisado por pares publicado no Congresso de Inicia\c{c}\~ao Cient\'ifica e Tecnol\'ogica (CICT); projeto PVF16985-2019, ``Gest\~ao Inteligente do Consumo de Energia'' (GERINCE).
\end{itemize}

% ------------------------------------------------------------------
% Projetos Open-Source Selecionados
% ------------------------------------------------------------------
\section{Projetos Open-Source Selecionados}

\project{FeaturesAnalyzer}{Python}{Toolkit de an\'alise de s\'eries temporais e machine learning para explora\c{c}\~ao de features, modelagem e previs\~ao sobre grandes volumes de dados.}
\project{ZSOM}{Python}{Implementa\c{c}\~ao do zero do algoritmo Self-Organizing Maps, base dos modelos n\~ao supervisionados da pesquisa de consumo de energia.}
\project{weasyprint-pdf-render}{Python, Docker}{Microsservi\c{c}o de renderiza\c{c}\~ao de HTML para PDF de alta performance, conteinerizado para escala horizontal.}
\project{linuwu-sense-dkms}{C, Shell}{Pacote de kernel DKMS para o driver Linuwu-Sense (Acer WMI Extras), demonstrando trabalho de baixo n\'ivel em Linux.}
\vspace{-3pt}

% ------------------------------------------------------------------
% Formacao
% ------------------------------------------------------------------
\section{Forma\c{c}\~ao Acad\^emica}
\textbf{UFRN -- Universidade Federal do Rio Grande do Norte} \hfill 2019 -- 2025 \\
Bacharelado em Sistemas de Informa\c{c}\~ao --- TCC com nota 10/10 \\[6pt]
\textbf{IMD/UFRN -- Instituto Metr\'opole Digital} \hfill 2018 -- 2019 \\
T\'ecnico em Redes de Computadores

% ------------------------------------------------------------------
% Idiomas & Premios
% ------------------------------------------------------------------
\section{Idiomas \& Pr\^emios}
\begin{itemize}[leftmargin=14pt, itemsep=2pt, topsep=0pt]
    \item \textbf{CodersRank:} 1\textsuperscript{o} lugar mundial em GraphQL (atualmente 2\textsuperscript{o}); 1\textsuperscript{o} em Jo\~ao Pessoa e Top 150 no Brasil.
    \item \textbf{Pr\^emios:} 2\textsuperscript{o} lugar --- 1\textsuperscript{a} Maratona de Programa\c{c}\~ao da UFRN (p\'olo regional); Bolsa de Pesquisa PVF16985-2019 (GERINCE).
    \item \textbf{Idiomas:} Portugu\^es (Nativo), Ingl\^es (Profissional), Espanhol (B\'asico -- Compreens\~ao).
\end{itemize}

\end{document}
